% from ega4-2-fr.tex
\oldpage[IV-2]{156}

\item[(ii)] L'hypothèse que $\mathcal J$ est localement nilpotent entraîne que
les espaces sous-jacents à $X$ et à $Y$ sont les mêmes, d'où les deux premières
assertions de (ii), compte tenu de (5.1.12.1); il reste à prouver (6.10.3.2),
car on en déduira aussitôt (6.10.3.3) par différence. Soit
$(f_i)_{1\leq i\leq p}$ une suite finie d'éléments de l'idéal maximal de
$\mathcal O_{X,x}$ dont les images dans
$\mathcal O_{Y,x}=\mathcal O_{X,x}/\mathcal J_x$ forment une suite
$\mathcal O_{Y,x}$-régulière maximale; l'hypothèse sur $\mathcal F$ et
$\mathcal J$ entraîne que le $\mathcal O_{Y,x}$-module
$\operatorname{gr}_{\mathcal J_x}^{\bullet}(\mathcal F_x)$ est libre de type
fini, donc la suite $(f_i)$ est
$\operatorname{gr}_{\mathcal J_x}^{\bullet}(\mathcal F_x)$-régulière; on en
déduit que cette suite est aussi $\mathcal F_x$-régulière (0, 15.1.19). Soit
d'autre part $n$ le plus grand entier tel que
$\mathcal F_x^{(n)}=\mathcal J_x^n\mathcal F_x\ne0$. L'hypothèse entraîne
aussi que
$\operatorname{gr}_{\mathcal J_x}^{\bullet}
(\mathcal F_x/\mathcal F_x^{(n)})$ est libre de type fini, donc la suite
$(f_i)$ est aussi $(\mathcal F_x/\mathcal F_x^{(n)})$-régulière
(\emph{loc. cit.}). Appliquant le lemme (3.4.1.4), on en conclut par récurrence
sur $i$, une suite exacte
\[
0\longrightarrow
\mathcal F_x^{(n)}/\left(\sum_{i=1}^{p}f_i\mathcal F_x^{(n)}\right)
\longrightarrow
\mathcal F_x/\left(\sum_{i=1}^{p}f_i\mathcal F_x\right).
\]
Mais l'hypothèse entraîne que $\mathcal F_x^{(n)}$ est un
$\mathcal O_{Y,x}$-module libre de type fini et $\ne0$, donc
$\mathcal F_x^{(n)}/(\sum_{i=1}^{p}f_i\mathcal F_x^{(n)})$ est isomorphe à
un module de la forme
$(\mathcal O_{Y,x}/(\sum_{i=1}^{p}f_i\mathcal O_{Y,x}))^m$ avec $m>0$; comme
\[
\operatorname{prof}\left(\mathcal O_{Y,x}/
\left(\sum_{i=1}^{p}f_i\mathcal O_{Y,x}\right)\right)=0
\]
(0, 16.4.6), on a aussi, par la caractérisation (0, 16.4.6, (i)) des modules
de profondeur nulle,
\[
\operatorname{prof}\left(\mathcal F_x^{(n)}/
\left(\sum_{i=1}^{p}f_i\mathcal F_x^{(n)}\right)\right)=0,
\]
puis
\[
\operatorname{prof}\left(\mathcal F_x/
\left(\sum_{i=1}^{p}f_i\mathcal F_x\right)\right)=0;
\]
les $(f_i)$ appartenant à l'idéal maximal de $\mathcal O_{Y,x}$ et formant
une suite $\mathcal F_x$-régulière, cela montre que l'on a
$\operatorname{prof}(\mathcal F_x)=p$ (0, 16.4.6, (ii)).
\end{enumerate}
\end{proof}

\begin{corollary}[6.10.4]
\label{IV.2.s6.10.4-fr}
Soit $U_{S_n}(\mathcal F)$ (resp. $U_{C_n}(\mathcal F)$) l'ensemble des
$x\in X$ tels que $\mathcal F$ vérifie $(S_n)$ au point $x$ (resp. tels que
$\operatorname{coprof}(\mathcal F_x)\leq n$). Si $\mathcal F$ est normalement
plat le long de $Y$, si $\operatorname{Supp}(\mathcal F)=X$ et si $\mathcal J$
est localement nilpotent, on a
$U_{S_n}(\mathcal F)=U_{S_n}(\mathcal O_Y)$ et
$U_{C_n}(\mathcal F)=U_{C_n}(\mathcal O_Y)$.
\end{corollary}

\begin{proposition}[6.10.5]
\label{IV.2.s6.10.5-fr}
Les notations étant celles de (6.10.1), supposons que $Y$ soit connexe et non
vide, et que $\mathcal F$ soit normalement plat le long de $Y$. Pour tout
entier $n\geq0$, on pose
\begin{equation}
r(n)=\operatorname{rg}_{\mathcal O_Y}
(\operatorname{gr}_{\mathcal J}^{n}(\mathcal F))
\tag{6.10.5.1}\label{IV.2.s6.10.5.1-fr}
\end{equation}
(le $\mathcal O_Y$-Module localement libre
$\operatorname{gr}_{\mathcal J}^{n}(\mathcal F)$ étant nécessairement de rang
constant). Alors:
\begin{enumerate}
\item[(i)] Il existe un polynôme $P\in\mathbf Q[T]$ tel que $P(n)=r(n)$ pour
tout $n$ assez grand.
\item[(ii)] On a $Y\cap\operatorname{Supp}(\mathcal F)=\varnothing$
(autrement dit $\mathcal F/\mathcal J\mathcal F=0$), ou
$Y\cap\operatorname{Supp}(\mathcal F)=Y$ (autrement dit
$Y\subset\operatorname{Supp}(\mathcal F)$).
\end{enumerate}
Dans le second cas, notons $d-1$ le degré de $P$; pour tout point maximal $y$
de $Y$, on a
\begin{equation}
\dim(\mathcal F_y)=d
\tag{6.10.5.2}\label{IV.2.s6.10.5.2-fr}
\end{equation}
et en particulier
\begin{equation}
\operatorname{codim}(Y,\operatorname{Supp}(\mathcal F))=d.
\tag{6.10.5.3}\label{IV.2.s6.10.5.3-fr}
\end{equation}
\begin{enumerate}
\item[(iii)] Supposons $Y\subset\operatorname{Supp}(\mathcal F)$. Pour tout
$x\in Y$, on a
\end{enumerate}
\begin{equation}
\dim(\mathcal F_x)=\dim(\mathcal O_{Y,x})+d.
\tag{6.10.5.4}\label{IV.2.s6.10.5.4-fr}
\end{equation}
De façon plus précise, il existe dans $\mathcal J_x$ des éléments $f_i$
$(1\leq i\leq d)$ faisant partie d'un système de paramètres pour
$\mathcal F_x$ (0, 16.3.6) et tels que, dans
$\operatorname{Spec}(\mathcal O_{X,x})$, on ait
\[
V(\mathcal J_x)=V\left(\sum_{i=1}^{d}f_i\mathcal O_{X,x}\right)
\cap\operatorname{Supp}(\mathcal F).
\]
\end{proposition}

\begin{proof}
Comme $Y$ est supposé connexe, la première assertion de (ii) découle de
(6.10.3, (i)). Si l'on a
$Y\cap\operatorname{Supp}(\mathcal F)=\varnothing$, l'assertion (i) est
triviale avec $P=0$. Supposons que
$\operatorname{Supp}(\mathcal F)\supset Y$, et soit $y$ un point maximal de
$Y$; comme on a alors $Y=\operatorname{Supp}(\mathcal F/\mathcal J\mathcal F)$,
$\mathcal F_y/\mathcal J_y\mathcal F_y$ est un
$\mathcal O_{X,y}$-module de longueur finie (3.1.4); posant
$s(n)=\operatorname{long}(\operatorname{gr}_{\mathcal J_y}^{n}(\mathcal F_y))$,
on sait qu'il y a un polynôme $R\in\mathbf Q[T]$ tel que $s(n)=R(n)$ pour
$n$ assez grand, savoir le polynôme $H(n+1)-H(n)$, où $H$ est le polynôme
de Hilbert--Samuel de $\mathcal F_y$ pour la filtration
$\mathcal J_y$-préadique (0, 16.2.1). Comme les
$\mathcal O_{Y,y}$-modules
$\operatorname{gr}_{\mathcal J_y}^{n}(\mathcal F_y)$ sont libres par
hypothèse, on a $r(n)=s(n)/m$, en désignant par $m$ la longueur du
$\mathcal O_{X,y}$-module $\mathcal O_{X,y}/\mathcal J_y$; on satisfait donc
à (i) en
prenant $P=R/m$. On sait en outre (0, 16.2.3) que
$\deg(H)=\dim(\mathcal F_y)$, d'où la relation (6.10.5.2); la relation
(6.10.5.3) s'en déduit à l'aide de (5.1.12.2).

